\section{Performance Fine-tuning}
\label{sec:performance-fine-tuning}
\postsec

The programs sampled by the program sampler have good coverage of the search space, but their qualities are not guaranteed.
This is because the optimization choices, such as tile structure and loop annotations, are all randomly sampled.
In this section, we introduce the performance tuner that fine-tunes the performance of the sampled programs via evolutionary search and a learned cost model.

The fine-tuning is performed iteratively. At each iteration, we first use evolutionary search to find a small batch of promising programs according to a learned cost model.  We then measure these programs on hardware to get the actual execution time cost. Finally, the profiling data got from measurement is used to re-train the cost model to make it more accurate.

The evolutionary search uses randomly sampled programs as well as high-quality programs from the previous measurement as the initial population and applies mutation and crossover to generate the next generation. The learned cost model is used to predict the \textit{fitness} of each program, which is the throughput of one program in our case.
We run evolution for a fixed number of generations and pick the best programs found during the search.
We leverage a learned cost model because the cost model can give relatively accurate estimations of the fitness of programs while being orders of magnitudes faster than the actual measurement.
It allows us to compare tens of thousands of programs in the search space in seconds, and pick the promising ones to do actual measurements.

\subsection{Evolutionary Search}
\label{subsec:evolutionary-search}

Evolutionary search~\cite{vikhar2016evolutionary} is a generic meta-heuristic algorithm inspired by biological evolution. By iteratively mutating high-quality programs, we can generate new programs with potentially higher quality. The evolution starts from the sampled initial generation.
To generate the next generation, we first select some programs from the current generation according to certain probabilities. The probability of selecting a program is proportional to its fitness predicted by the learned cost model (\autoref{subsec:learned-cost-model}), meaning that the program with a higher performance score has a higher probability to be selected.
For the selected programs, we randomly apply one of the evolution operations to generate a new program. Basically, for decisions we made during sampling (\autoref{subsec:random-sampling}), we design corresponding evolution operations to rewrite and fine-tune them.

\textbf{Tile size mutation.}
This operation scans the program and randomly selects a tiled loop.
For this tiled loop, it divides a tile size of one tile level by a random factor and multiplies this factor to another level.
Since this operation keeps the product of tile sizes equal to the original loop length, the mutated program is always valid.

\textbf{Parallel mutation.}
This operation scans the program and randomly selects a loop that has been annotated with parallel. 
For this loop, this operation changes the parallel granularity by either fusing its adjacent loop levels or splitting it by a factor.

\textbf{Pragma mutation.}
Some optimizations in a program are specified by compiler-specific pragma. This operation scans the program and randomly selects a pragma. For this pragma, this operation randomly mutates it into another valid value. For example, our underlying code generator supports auto unrolling with a maximum number of steps by providing an \texttt{auto\_unroll\_max\_step=N} pragma. We randomly tweak the number $N$.

\textbf{Computation location mutation.}
This operation scans the program and randomly selects a flexible node that is not multi-level tiled (\textit{e.g.}, a padding node in the convolution layer). For this node, the operation randomly changes its computation location to another valid attach point.

\textbf{Node-based crossover.}
Crossover is an operation to generate new offspring by combining the genes from two or more parents.
The genes of a program in \Ansor are its rewriting steps.
Every program generated by \Ansor is rewritten from its initial naive implementation.
\Ansor preserves a complete rewriting history for each program during sketch generation and random annotation.
We can treat rewriting steps as the genes of a program because they describe how this program is formed from the initial naive one.
Based on this, we can generate a new program by combining the rewriting steps of two existing programs.
However, arbitrarily combining rewriting steps from two programs might break the dependencies in steps and create an invalid program.
As a result, the granularity of crossover operation in \Ansor is based on nodes in the DAG, because the rewriting steps across different nodes usually have less dependency.
Ansor randomly selects one parent for each node and merges the rewriting steps of selected nodes.
When there are dependencies between nodes, \Ansor tries to analyze and adjust the steps with simple heuristics. \Ansor further verifies the merged programs to guarantee the functional correctness. The verification is simple because \Ansor only uses a small set of loop transformation rewriting steps, and the underlying code generator can check the correctness by dependency analysis.

The evolutionary search leverages mutation and crossover to generate a new set of candidates repeatedly for several rounds and outputs a small set of programs with the highest scores.
These programs will be compiled and measured on the target hardware to obtain the real running time cost.
The collected measurement data is then used to update the cost model.
In this way, the accuracy of the learned cost model is gradually improved to match the target hardware.
Consequently, the evolutionary search gradually generates higher-quality programs for the target hardware platform.

Unlike the search algorithms in TVM and FlexTensor that can only work in a fixed grid-like parameter space, the evolutionary operations in \Ansor are specifically designed for tensor programs. They can be applied to general tensor programs and can handle a search space with complicated dependency.
Unlike the unfolding rules in Halide auto-scheduler, these operations can perform out-of-order modifications to programs, addressing the sequential limitations.

\subsection{Learned Cost Model}
\label{subsec:learned-cost-model}

A cost model is necessary for estimating the performance of programs quickly during the search. We adopt a learned cost model similar to related works \cite{adams2019learning, chen2018learning} with newly designed program features.
A system based on learned cost models has great portability because a single model design can be reused for different hardware backends by feeding in different training data.

Since our target programs are mainly data parallel tensor programs, which are made by multiple interleaved loop nests with several assignment statements as the innermost statements, we train the cost model to predict the score of one innermost non-loop statement in a loop nest.
For a full program, we make predictions for each innermost non-loop statement and add the predictions up as the score. We build the feature vector for an innermost non-loop statement by extracting features in the context of a full program. The extracted features include arithmetic features and memory access features.
A detailed list of extracted features is in \autoref{sec:extracted-features}.

We use weighted squared error as the loss function. Because we mostly care about identifying the well-performing programs from the search space, we put more weight on the programs that run faster.
Specifically, the loss function of the model $f$ on a program $P$ with throughput $y$ is 
$ loss(f, P, y) = w_p (\sum_{s \in S(P)} f(s) - y )^2 =  y (\sum_{s \in S(P)} f(s) - y)^2  $
where $S(P)$ is the set of innermost non-loop statements in $P$. We directly use the throughput $y$ as weight.
We train a gradient boosting decision tree \cite{chen2016xgboost} as the underlying model $f$. A single model is trained for all tensor programs coming from all DAGs, and we normalize the throughput of all programs coming from the same DAG to be in the range of $[0, 1]$.
When optimizing a DNN, the number of measured programs are typically less than 30,000.
Training a gradient boosting decision tree is very fast on such a small data sets, so we train a new model every time instead of doing incremental updates.






